\documentclass[12pt,a4paper]{article}
\usepackage[utf8]{inputenc}
\usepackage[T1]{fontenc}
\usepackage{amsmath,amssymb,amsthm}
\usepackage[margin=1in]{geometry}
\usepackage{hyperref}
\usepackage{enumitem}
\usepackage{array}

\newtheorem{theorem}{Theorem}
\newtheorem{lemma}[theorem]{Lemma}
\newtheorem{proposition}[theorem]{Proposition}
\newtheorem{corollary}[theorem]{Corollary}
\newtheorem{remark}[theorem]{Remark}
\newtheorem{definition}[theorem]{Definition}

\newcommand{\so}{\mathfrak{so}}
\newcommand{\Ad}{\operatorname{Ad}}
\newcommand{\Lie}{\operatorname{Lie}}
\newcommand{\GL}{\operatorname{GL}}
\newcommand{\AGL}{\operatorname{AGL}}
\newcommand{\PGL}{\operatorname{PGL}}
\newcommand{\SL}{\operatorname{SL}}
\newcommand{\GF}{\operatorname{GF}}

\title{Quantum-Mimetic Ternary Systems on $\mathbb{R}^{27}$}

\author{ATroubledSnake\\\texttt{atroubledsnake@pm.me}}
\date{}

\begin{document}
\maketitle

\begin{abstract}
We analyze the linear-algebraic structure of a classical analog system
consisting of three balanced ternary wires, each carrying a real amplitude
in $\mathbb{R}^3$, combined via tensor product into $\mathbb{R}^{27}$.
We prove that local $\mathfrak{so}(3)$ generators on each wire together
with the ternary CNOT gate on any two wires sharing a common wire generate
the full Lie algebra $\mathfrak{so}(27)$ (dimension 351). In contrast,
a single CNOT generates only $\mathfrak{so}(9) \oplus \mathfrak{so}(3)$
(dimension 39). The 27 computational basis states form an orthonormal
spherical code on $S^{26}$ achieving 50\% of the Rankin bound. We define
a winner-take-all measurement rule and analyze its noise stability.
The discrete symmetry group generated by CNOT gates and local involutions
is isomorphic to $\GL(3,3)$ (order $11{,}232$), acting transitively on the
26 non-zero basis states ($|000\rangle$ is fixed). These results provide the minimal mathematical structure
required to describe a ``quantum-mimetic'' classical system whose gate
set is universal in $O(27)$.
\end{abstract}

\section{Introduction}

\subsection{Motivation}

A balanced ternary digit (trit) takes values in $\{-1,0,1\}$ and carries
$\log_2 3 \approx 1.585$ bits of information. The idea of using the tensor
product of three ternary wires to produce a $27$-dimensional state space
raises a natural question: can linear gates acting locally on each wire,
together with a single coupling gate, generate the full orthogonal group
$O(27)$? Or does the system remain confined to a proper submanifold
equivalent to three independent wires?

Ternary computing has a long and storied history. The Setun computer
(Brusentsov, Moscow State University, 1958--1965) was the first and only
ternary computer built at scale, using balanced ternary digits implemented
with magnetic amplifiers~\cite{brusentsov1965}. The advantages of balanced
ternary over binary include a simpler arithmetic (no separate sign bit,
rounding = truncation), a denser information representation ($\log_2 3$
vs.\ 1 bit per wire), and symmetric logic around zero. Despite these
advantages, the dominance of binary CMOS ultimately marginalized ternary
logic at the hardware level.

However, the recent development of metaplectic anyons in topological
quantum computing~\cite{bocharov2016improved} has revived interest in
ternary quantum gates. In the metaplectic anyon model based on $SU(2)_4$,
the natural computational basis is ternary (qutrit-based), and controlled
gates analogous to CNOT arise as braiding operations. This has led to a
substantial body of work on ternary quantum arithmetic, including ripple-carry
and carry-lookahead adders~\cite{bocharov2016improved}, multiplier
circuits~\cite{ilyas2022ternary}, and optimized ternary gate
synthesis~\cite{koleTernary}.

This paper takes a different perspective. Rather than building a quantum
computer from ternary units, we consider a \emph{classical} analog system
whose state space is the tensor product of three real vector spaces
$\mathbb{R}^3$ (one per ternary wire), yielding $\mathbb{R}^{27}$. The
``gates'' are linear orthogonal transformations on $\mathbb{R}^{27}$
generated by local $\mathfrak{so}(3)$ rotations on each wire together
with the ternary CNOT coupling gate. We ask: what is the reachable set
of states from a given initial state? Is the system universal in
$O(27)$, or does it remain confined to a proper submanifold?

\subsection{Summary of Results}

Our main results are:

\begin{enumerate}
\item \textbf{Two-wire universality (Theorem~\ref{thm:2wire}).} The Lie
algebra generated by local $\mathfrak{so}(3)$ generators on two wires and
the ternary CNOT gate is $\mathfrak{so}(9)$ (dimension 36). The proof uses
Dynkin's classification of maximal subalgebras: $\mathfrak{so}(3)\oplus
\mathfrak{so}(3)$ is a Type~S maximal subalgebra of $\mathfrak{so}(9)$,
and the CNOT gate lies outside the corresponding subgroup
$SO(3)\times SO(3)$.

\item \textbf{Three-wire universality (Theorem~\ref{thm:3wire}).} With
three wires, two CNOT gates sharing a common wire generate the full
$\mathfrak{so}(27)$ (dimension 351). One CNOT generates only
$\mathfrak{so}(9)\oplus\mathfrak{so}(3)$ (dimension 39).

\item \textbf{Spanning tree generation (Theorem~\ref{thm:ktree}).} For
$k\ge 2$ wires, CNOT gates on edges forming a spanning tree generate
$\mathfrak{so}(3^k)$. The minimal number of CNOT gates required is $k-1$.

\item \textbf{Discrete symmetry group (Proposition~\ref{prop:gl33}).}
The discrete group generated by CNOT gates (both directions) on all wire
pairs together with the local inversion gates $T_{\mathrm{inv}}$ on each
wire is isomorphic to $\GL(3,3)$ (the general linear group over
$\GF(3)^3$), of order $11{,}232$, acting transitively on the 26 non-zero
basis states ($|000\rangle$ is fixed).

\item \textbf{Spherical code.} The 27 computational basis states form a
Parseval frame on $S^{26}$, achieving 50\% of the Rankin bound
(27 vs.\ 54 antipodal points). The winner-take-all measurement has a noise
threshold $\sigma < 0.34$ for reliable communication.
\end{enumerate}

\subsection{Relation to Prior Work}

The study of ternary quantum gates in the metaplectic anyon framework
was pioneered by Cui and Wang~\cite{cuietal2015}, who showed that the
controlled-$X$ gate (the ternary CNOT analog) together with single-qutrit
Clifford gates generates the full ternary Clifford group. Bocharov
et al.~\cite{bocharov2016improved} extended this to arithmetic circuits
and proved that Clifford + $C(X)$ generates $\GL(n, \mathbb{F}_3) \ltimes
\mathbb{F}_3^n$ as a maximal subgroup of the symmetric group. Our work
complements these results from the Lie algebra perspective: we analyze
the \emph{continuous} orthogonal transformations generated by the same
gates, showing that the reachable space is the full sphere $S^{26}$.

In classical ternary computing, the state space has typically been modeled
as a finite set of discrete levels (e.g., $\{-V,0,+V\}$ in voltage-mode
logic). Our analog model differs fundamentally: the amplitudes are
continuous real numbers, and the computational basis states are just 27
special points on a continuous sphere. This is analogous to the
relationship between qubit-based quantum computing (continuous
amplitudes in $\mathbb{C}^{2^n}$) and classical digital logic (discrete
$0/1$ states). The present work provides the minimal mathematical
framework for a ``quantum-mimetic'' analog ternary computer.

\section{State Space}

Let $\mathcal{V} = \mathbb{R}^3$ with ordered orthonormal basis
\[
\mathcal{B} = \{|-1\rangle, |0\rangle, |1\rangle\},\qquad
|-1\rangle = \begin{pmatrix}1\\0\\0\end{pmatrix},\;
|0\rangle = \begin{pmatrix}0\\1\\0\end{pmatrix},\;
|1\rangle = \begin{pmatrix}0\\0\\1\end{pmatrix}.
\]

Three wires $\mathcal{A}, \mathcal{B}, \mathcal{C}$ produce the tensor
product space
\[
\mathcal{H} = \mathcal{V}_{\mathcal{A}} \otimes \mathcal{V}_{\mathcal{B}}
\otimes \mathcal{V}_{\mathcal{C}} \cong \mathbb{R}^{27}.
\]

A macro-state is
\[
|\Psi\rangle = \sum_{i,j,k \in \{-1,0,1\}} c_{ijk}\,
|i\rangle_{\mathcal{A}} \otimes |j\rangle_{\mathcal{B}} \otimes
|k\rangle_{\mathcal{C}},\qquad \sum_{i,j,k} c_{ijk}^2 = 1.
\]

The 27 basis vectors are orthonormal and lie on the unit sphere $S^{26}$.

\section{Lie Algebra Generation}

\subsection{Local Generators}

The Lie algebra $\mathfrak{so}(3)$ has basis
\[
L_x = \begin{pmatrix}0&0&0\\0&0&-1\\0&1&0\end{pmatrix},\;
L_y = \begin{pmatrix}0&0&1\\0&0&0\\-1&0&0\end{pmatrix},\;
L_z = \begin{pmatrix}0&-1&0\\1&0&0\\0&0&0\end{pmatrix},
\]
satisfying $[L_x, L_y] = L_z$ and cyclic permutations.

These are lifted to $\mathfrak{so}(27)$ via Kronecker products:
\[
L_{A}(X) = X \otimes I \otimes I,\quad
L_{B}(X) = I \otimes X \otimes I,\quad
L_{C}(X) = I \otimes I \otimes X.
\]

For two wires, we write
\[
L_{A}(X) = X \otimes I,\qquad
L_{B}(X) = I \otimes X,
\]
acting on $\mathbb{R}^9 = \mathbb{R}^3 \otimes \mathbb{R}^3$.

The 9 local generators generate $\mathfrak{so}(3) \oplus \mathfrak{so}(3)
\oplus \mathfrak{so}(3)$, a 9-dimensional proper subalgebra of
$\mathfrak{so}(27)$.

\subsection{The Ternary CNOT Gate}

The ternary CNOT acts on two wires as
\[
\text{CNOT}: |i,j\rangle \mapsto |i,\, j \oplus i\rangle,
\]
where $\oplus$ denotes addition modulo 3. As a $9 \times 9$ permutation
matrix, it lies in $SO(9)$ (determinant 1; all eigenvalues are third roots
of unity). Explicitly, ordering the basis as
$|00\rangle,|01\rangle,|02\rangle,|10\rangle,|11\rangle,|12\rangle,
|20\rangle,|21\rangle,|22\rangle$, the CNOT matrix is
\[
\text{CNOT} = \begin{pmatrix}
I & 0 & 0 \\
0 & S & 0 \\
0 & 0 & S^2
\end{pmatrix},
\quad
S = \begin{pmatrix}0&0&1\\1&0&0\\0&1&0\end{pmatrix},
\]
where $S$ is the cyclic shift on $\mathbb{R}^3$ ($S^3 = I$). The CNOT is
an orthogonal matrix with $\det(\text{CNOT}) = 1$.

Since all eigenvalues are third roots of unity ($1,\omega,\omega^2$ with
$\omega = e^{2\pi i/3}$), its matrix logarithm
$G_{AB} = \log(\text{CNOT})$ exists as a real skew-symmetric matrix in
$\mathfrak{so}(9)$. Explicitly, the eigenvalues of $\log(\text{CNOT})$
are $0$ (multiplicity 5) and $\pm 2\pi i/3$ (multiplicity 2 each), all
purely imaginary, confirming $G_{AB} \in \mathfrak{so}(9)$.

Lifted to $\mathbb{R}^{27}$, the CNOT on wires $A$ and $B$ with wire $C$
as spectator is
\[
G_{AB}^{(27)} = \log(\text{CNOT}) \otimes I \in \mathfrak{so}(27).
\]

\subsection{Generation Theorem: The Correct Formulation}
\label{sec:groupformulation}

Let $G$ be the group generated by the local unitary groups
$\exp(\mathfrak{so}(3)_A)$, $\exp(\mathfrak{so}(3)_B)$, and the CNOT gate.
Formally,
\[
G = \langle \exp(\mathfrak{so}(3)_A),\,
\exp(\mathfrak{so}(3)_B),\, \text{CNOT} \rangle \subset SO(9).
\]

Let $L = \Lie(G)$ be the Lie algebra of $G$. Since $\text{CNOT}$ is a
discrete element, its adjoint action $\Ad_{\text{CNOT}}$ is an automorphism
of $L$. Therefore $L$ is the smallest Lie subalgebra of $\mathfrak{so}(9)$
that contains $\mathfrak{so}(3)_A \oplus \mathfrak{so}(3)_B$ and is
invariant under $\Ad_{\text{CNOT}}$. Equivalently,
\[
L = \big\langle \mathfrak{so}(3)_A \oplus \mathfrak{so}(3)_B,\,
\Ad_{\text{CNOT}}(\mathfrak{so}(3)_A \oplus \mathfrak{so}(3)_B),\,
\Ad_{\text{CNOT}}^2(\mathfrak{so}(3)_A \oplus \mathfrak{so}(3)_B),
\ldots \big\rangle_{\text{Lie}},
\]
where $\langle \cdot \rangle_{\text{Lie}}$ denotes Lie closure (repeated
commutators). Since $\text{CNOT}^3 = I$, the adjoint action has order 3,
so only $\Ad_{\text{CNOT}}$ and $\Ad_{\text{CNOT}}^2$ need be considered.

\subsection{Maximal Subalgebra Argument}

The key to the proof is the following result from Dynkin's classification
of maximal subalgebras of classical Lie algebras.

\begin{lemma}[Dynkin classification]
\label{lemma:dynkin}
Let $m,n \ge 3$ with $(m,n) \neq (4,4)$. The tensor product embedding
\[
\mathfrak{so}(m) \oplus \mathfrak{so}(n) \to \mathfrak{so}(mn)
\]
via the irreducible representation $\mathbb{R}^m \otimes \mathbb{R}^n$
is a maximal subalgebra of $\mathfrak{so}(mn)$. In Dynkin's
classification~\cite{dynkin1952}, this is called a \emph{non-regular
non-simple} subalgebra (Type~III); it belongs to the broader class of
Type~S (or $S$-subalgebras) in Dynkin's terminology
\cite{dynkin1952,bocharov2016improved}.
The exceptional cases $m,n \in \{2,4,8\}$ are excluded because
$\mathfrak{so}(4) \cong \mathfrak{so}(3)\oplus\mathfrak{so}(3)$ is not
simple, $\mathfrak{so}(2)$ is abelian, and $\mathfrak{so}(8)$ has triality
(an $S_3$ outer automorphism group). The case $m=n=4$ is excluded because
$\mathfrak{so}(4)\oplus\mathfrak{so}(4)$ fails to be maximal in
$\mathfrak{so}(16)$ due to the decomposition
$\bigwedge^2(\mathbb{R}^4\otimes\mathbb{R}^4)$ containing additional
invariants under the $\mathfrak{so}(3)\oplus\mathfrak{so}(3)$ factors.
\end{lemma}

\begin{proof}
See Dynkin~\cite[\S15, no.~46, Type~III table (p.~114)]{dynkin1952},
where the non-regular non-simple maximal subalgebras of the classical Lie
algebras are tabulated. For $\mathfrak{so}(N)$, the Type~III maximal
subalgebras include
\[
\mathfrak{so}(s) \oplus \mathfrak{so}(t) \subset \mathfrak{so}(st)
\qquad(3 \le s \le t,\; s,t \neq 4).
\]
The representation $\mathbb{R}^m \otimes \mathbb{R}^n$ of
$\mathfrak{so}(m)\oplus\mathfrak{so}(n)$ is irreducible, and the
dimension count
\[
\dim(\mathfrak{so}(m)\oplus\mathfrak{so}(n)) = \frac{m(m-1)}{2} +
\frac{n(n-1)}{2}
\]
is strictly less than $\dim\mathfrak{so}(mn) = \frac{mn(mn-1)}{2}$.
Maximality fails only when the representation is not faithful (e.g.,
$\mathfrak{so}(2)$ gives a 1-dimensional center) or when the algebra is
not simple ($\mathfrak{so}(4)$) or when special outer automorphisms
exist ($\mathfrak{so}(8)$ triality).
\end{proof}

\begin{corollary}
\label{cor:so3so3max}
The subalgebra $\mathfrak{so}(3)_A \oplus \mathfrak{so}(3)_B$ is a maximal
subalgebra of $\mathfrak{so}(9)$. It arises from the Type~S tensor product
embedding with $m=n=3$.
\end{corollary}

\subsection{Two-Wire Generation}

\begin{theorem}[Two-wire generation]
\label{thm:2wire}
For $k=2$, the Lie algebra generated by $\mathfrak{so}(3)_A$,
$\mathfrak{so}(3)_B$, and $\text{CNOT}(A,B)$ (in the sense of
Section~\ref{sec:groupformulation}) is the full $\mathfrak{so}(9)$
(dimension 36).
\end{theorem}

\begin{proof}
Let $L$ be the Lie algebra generated by $\mathfrak{so}(3)_A \oplus
\mathfrak{so}(3)_B$ and CNOT, i.e., $L = \Lie(G)$ where
$G = \langle \exp(\mathfrak{so}(3)_A \oplus \mathfrak{so}(3)_B),
\text{CNOT} \rangle$.

First, we show that $\text{CNOT} \notin SO(3) \times SO(3)$. The subgroup
$SO(3) \times SO(3)$ consists of matrices of the form $R_A \otimes R_B$
with $R_A,R_B \in SO(3)$. Its Lie algebra is $\mathfrak{so}(3)_A \oplus
\mathfrak{so}(3)_B$, whose elements are $X \oplus Y = X \otimes I + I \otimes Y$
for $X,Y \in \mathfrak{so}(3)$.

Consider the matrix logarithm $G = \log(\text{CNOT}) \in \mathfrak{so}(9)$.
Its eigenvalues are $0$ (multiplicity 5) and $\pm 2\pi i/3$ (multiplicity 2
each). Now let $X \in \mathfrak{so}(3)$ have eigenvalues $\{0, \pm i\theta\}$
for some $\theta$. Then $X \otimes I$ has eigenvalues $\{0, \pm i\theta\}$
each with multiplicity 3; similarly $I \otimes Y$ has eigenvalues
$\{0, \pm i\phi\}$ each with multiplicity 3. Their sum $X \otimes I + I \otimes Y$
has eigenvalues that are sums of pairs of eigenvalues of $X$ and $Y$, with
multiplicities given by products of the individual multiplicities. All
eigenvalues of $X \otimes I + I \otimes Y$ therefore occur with
multiplicities that are multiples of 3 (since each eigenvalue class of $X$
or $Y$ has multiplicity 3). In particular, the multiplicity of the zero
eigenvalue in any element of $\mathfrak{so}(3)_A \oplus \mathfrak{so}(3)_B$
is divisible by 3 (the zero eigenvalue of $X$ paired with the zero eigenvalue
of $Y$ gives multiplicity $3 \times 3 = 9$; other zero combinations add
multiples of 3).

The eigenvalue $0$ of $G$ has multiplicity $5$, which is not divisible by 3.
Therefore $G \notin \mathfrak{so}(3)_A \oplus \mathfrak{so}(3)_B$, and
consequently $\text{CNOT} \notin SO(3) \times SO(3)$.

Since $\text{CNOT} \notin SO(3) \times SO(3)$, the group $G$ is strictly
larger than $SO(3) \times SO(3)$. Therefore $L$ strictly contains
$\mathfrak{so}(3)_A \oplus \mathfrak{so}(3)_B$.

By Corollary~\ref{cor:so3so3max}, $\mathfrak{so}(3)_A \oplus
\mathfrak{so}(3)_B$ is a maximal subalgebra of $\mathfrak{so}(9)$.
Since $L$ strictly contains a maximal subalgebra, we must have
$L = \mathfrak{so}(9)$.
\end{proof}

\subsection{Explicit Commutator Verification}

The maximal subalgebra argument above is sufficient for the proof, but
we can also verify directly that the CNOT generator couples the two wires.
Let
\[
J_B^z = I \otimes L_z \in \mathfrak{so}(3)_B,\qquad
J_B^x = I \otimes L_x \in \mathfrak{so}(3)_B,
\]
and consider the adjoint action of CNOT on $J_B^z$. Writing CNOT in
block form as $\operatorname{diag}(I, S, S^2)$ where $S$ is the cyclic
shift, we obtain
\[
\Ad_{\text{CNOT}}(J_B^z) = \text{CNOT} \cdot (I \otimes L_z) \cdot
\text{CNOT}^{-1}
= \begin{pmatrix} L_z & 0 & 0 \\ 0 & S L_z S^{-1} & 0 \\
0 & 0 & S^2 L_z S \end{pmatrix}.
\]
A brief computation gives $S L_z S^{-1} = L_y$ and
$S^2 L_z S = L_x$, so this becomes
$\operatorname{diag}(L_z, L_y, L_x)$ in the $3 \times 3$ block
decomposition.

This element is not a pure tensor $I \otimes (\cdot)$: it has different
$3\times 3$ blocks on each diagonal sector, mixing the $A$ and $B$
degrees of freedom. A brief computation confirms
\[
[J_B^z,\; \Ad_{\text{CNOT}}(J_B^z)]
\]
has non-zero projection onto the orthogonal complement of
$\mathfrak{so}(3)_A \oplus \mathfrak{so}(3)_B$, confirming that
$\mathfrak{so}(3)_A \oplus \mathfrak{so}(3)_B$ is not an ideal in $L$
and that repeated commutation fills $\mathfrak{so}(9)$.

\begin{theorem}[Three-wire generation]
\label{thm:3wire}
For $k=3$, the Lie algebra generated by $\mathfrak{so}(3)_A$,
$\mathfrak{so}(3)_B$, $\mathfrak{so}(3)_C$, and any two CNOT gates
sharing a common wire (e.g., $\text{CNOT}(A,B)$ and $\text{CNOT}(A,C)$)
is the full $\mathfrak{so}(27)$ (dimension 351). A single CNOT generates
only $\mathfrak{so}(9) \oplus \mathfrak{so}(3)$ (dimension 39).
\end{theorem}

\begin{proof}
The single-CNOT case follows from Theorem~\ref{thm:2wire} on wires
$A,B$ together with the disjoint-support lemma: a generator on wire $C$
commutes with all elements acting only on $A,B$, so the algebra
decomposes as $\mathfrak{so}(9)_{AB} \oplus \mathfrak{so}(3)_C$.

For the two-CNOT case, let $L$ be the generated algebra (in the sense of
Section~\ref{sec:groupformulation}). By Theorem~\ref{thm:2wire} applied
to $(A,B)$ and $(A,C)$, we have $\mathfrak{so}(9)_{AB} \subset L$ and
$\mathfrak{so}(9)_{AC} \subset L$. Both contain $\mathfrak{so}(3)_A$
(the local algebra on wire $A$).

Now consider the subalgebra $\mathfrak{so}(9)_{AB} \oplus \mathfrak{so}(3)_C
\subset \mathfrak{so}(27)$. By Lemma~\ref{lemma:dynkin} with $m=9$, $n=3$,
the embedding $\mathfrak{so}(9)\oplus\mathfrak{so}(3) \to
\mathfrak{so}(27)$ via $\mathbb{R}^9\otimes\mathbb{R}^3$ is a Type~S maximal
subalgebra of $\mathfrak{so}(27)$ (both $9$ and $3$ are $\neq 2,4,8$).

We claim that $\mathfrak{so}(9)_{AC}$ is not contained in
$\mathfrak{so}(9)_{AB}\oplus\mathfrak{so}(3)_C$. To see this explicitly,
consider the generator
\[
G = \Ad_{\text{CNOT}_{AC}}(J_C^z) = \text{CNOT}_{AC} \cdot (I \otimes I
\otimes L_z) \cdot \text{CNOT}_{AC}^{-1} \in \mathfrak{so}(9)_{AC}.
\]
In the $3\times 3$ block decomposition relative to wire $A$ (as in
Section~\ref{sec:explicitcomm}),
\[
G = \operatorname{diag}(L_z, L_y, L_x)
\]
acting on wires $A$ and $C$ with wire $B$ as spectator; equivalently
\[
G = \sum_{i=0}^{2} |i\rangle\langle i|_A \otimes I_B \otimes M_i,
\qquad M_0 = L_z,\; M_1 = L_y,\; M_2 = L_x.
\]

Now suppose $G \in \mathfrak{so}(9)_{AB}\oplus\mathfrak{so}(3)_C$. Any
element of $\mathfrak{so}(9)_{AB}\oplus\mathfrak{so}(3)_C$ has the form
$H \otimes I_C + I_{AB} \otimes K$ with $H \in \mathfrak{so}(9)$,
$K \in \mathfrak{so}(3)$. Projecting onto the $|i\rangle\langle i|_A$
sector by multiplying on both sides by $|i\rangle\langle i|_A \otimes I_B$,
we obtain
\[
|i\rangle\langle i|_A \otimes I_B \otimes M_i
= |i\rangle\langle i|_A \otimes X_i \otimes I_C
+ |i\rangle\langle i|_A \otimes I_B \otimes K,
\]
where $X_i = (|i\rangle\langle i|_A \otimes I_B)\, H\,
(|i\rangle\langle i|_A \otimes I_B) \in \mathfrak{so}(3)_B$ is the
$B$-only component remaining after the $A$-projection. Comparing both
sides,
\[
X_i \otimes I_C + I_B \otimes K = I_B \otimes M_i,\qquad i=0,1,2.
\]
Taking the partial trace over wire $C$ of each equation gives
$3\,X_i = 0$, hence $X_i = 0$ for all $i$. Substituting back,
$I_B \otimes K = I_B \otimes M_i$, so $K = M_i$ simultaneously for
$i=0,1,2$, which is impossible since $L_z$, $L_y$, $L_x$ are distinct.
Therefore $G \notin \mathfrak{so}(9)_{AB} \oplus \mathfrak{so}(3)_C$, and
consequently $\mathfrak{so}(9)_{AC} \not\subset
\mathfrak{so}(9)_{AB}\oplus\mathfrak{so}(3)_C$.

Since $L$ contains $\mathfrak{so}(9)_{AB}\oplus\mathfrak{so}(3)_C$
(because $\mathfrak{so}(9)_{AB} \subset L$ and $\mathfrak{so}(3)_C
\subset L$) and also contains $\mathfrak{so}(9)_{AC}$ which is not
contained in this subalgebra, $L$ strictly contains the maximal
subalgebra. Therefore $L = \mathfrak{so}(27)$.

Explicitly, one can verify that the commutator of elements from
$\mathfrak{so}(9)_{AB}$ and $\mathfrak{so}(9)_{AC}$ generates components
with support on all three wires. Let $Y = \Ad_{\text{CNOT}_{AC}}(J^z_C)
\in \mathfrak{so}(9)_{AC}$. Choose $X = [J^x_A,
\Ad_{\text{CNOT}_{AB}}(J^z_B)] \in \mathfrak{so}(9)_{AB}$. Writing
$\Ad_{\text{CNOT}_{AB}}(J^z_B) = \sum_{a} |a\rangle\langle a|_A \otimes
X_a \otimes I_C$ with $X_a = S^a L_z S^{-a}$, and $J^x_A = L_x \otimes
I_B \otimes I_C$, we obtain
\[
X = \sum_{p,q} (L_x)_{pq}\, |p\rangle\langle q|_A \otimes (X_q - X_p)
\otimes I_C.
\]
Since $L_x$ has non-zero entries only at $(1,2)$ and $(2,1)$,
this expression is non-zero. The bracket $[X,Y]$ has components acting
on all three wires; its partial trace over $C$ vanishes while the
full commutator is non-zero, so it cannot lie in $\mathfrak{so}(9)_{AB}
\oplus \mathfrak{so}(3)_C$. (A direct numerical check confirms the
$27\times 27$ matrix $[X,Y]$ has norm $4\sqrt{2}$ and is orthogonal
to the direct sum.)
\end{proof}

\begin{remark}
The SWAP gate does not enlarge the algebra: $\text{SWAP}(A,B)$ is
contained in the Lie closure of the local generators, acting as a
discrete Weyl group element.
\end{remark}

\subsection{Disjoint Support Lemma}
\label{sec:disjoint}

\begin{lemma}[Disjoint support]
\label{lemma:disjoint}
Let $V = V_1 \otimes V_2$ be a tensor product of finite-dimensional
real vector spaces. For any $X \in \mathfrak{so}(V_1)$ and
$Y \in \mathfrak{so}(V_2)$, the lifted generators satisfy
\[
[X \otimes I_2,\; I_1 \otimes Y] = 0
\]
in $\mathfrak{so}(V_1 \otimes V_2)$. More generally, if
$A \in \mathfrak{so}(V_1 \otimes \cdots \otimes V_k)$ acts
non-trivially only on wires $S \subseteq \{1,\dots,k\}$ and
$B$ acts non-trivially only on $T \subseteq \{1,\dots,k\}$ with
$S \cap T = \varnothing$, then $[A,B] = 0$.
\end{lemma}

\begin{proof}
Immediate from the definition of the Kronecker product:
\[
(X \otimes I)(I \otimes Y) = X \otimes Y = (I \otimes Y)(X \otimes I).
\]
The general case follows by embedding each operator into the full tensor
product via $A \otimes I_{[k]\setminus S}$ and noting that disjoint
supports give commuting factors.
\end{proof}

\subsection{Generation Theorem for $k$ Wires}

\begin{theorem}[Spanning tree generation]
\label{thm:ktree}
For $k \ge 2$, the Lie algebra generated by $\mathfrak{so}(3)_i$ for
each wire $i = 0,\dots,k-1$ together with $\text{CNOT}_{i,j}$ for edges
forming a spanning tree on $k$ wires (exactly $k-1$ edges) is the full
$\mathfrak{so}(3^k)$.
\end{theorem}

\begin{proof}
We prove by induction on $k$. The base case $k=2$ is
Theorem~\ref{thm:2wire}.

Assume the theorem holds for $k-1$ wires. Let $T$ be a spanning tree on
$k$ wires with leaf $l$ adjacent to $p$. Let $V = \mathbb{R}^{3^{k-1}}$
(the state space of the $k-1$ wires excluding $l$) and
$W = \mathbb{R}^3$ (wire $l$).

By the induction hypothesis applied to the tree $T\setminus\{l\}$,
the Lie algebra $\mathfrak{so}(V)$ is contained in the generated
algebra $L$. By Theorem~\ref{thm:2wire} applied to wires $p$ and $l$,
$\mathfrak{so}(9)_{p,l} \subset L$ (in the lifted sense, with all other
wires as spectators). The intersection
$\mathfrak{so}(V) \cap \mathfrak{so}(9)_{p,l}$ contains
$\mathfrak{so}(3)_p$ (the local algebra on wire $p$), so both
subalgebras are non-trivially coupled.

Now consider the subalgebra $\mathfrak{so}(V) \oplus \mathfrak{so}(W)
\subset \mathfrak{so}(V\otimes W) = \mathfrak{so}(3^k)$. The embedding
$\mathfrak{so}(V)\oplus\mathfrak{so}(W) \to \mathfrak{so}(V\otimes W)$
via the tensor product $\mathbb{R}^{3^{k-1}}\otimes\mathbb{R}^3$ is a
Type~S maximal subalgebra by Lemma~\ref{lemma:dynkin}: here $m = 3^{k-1}$
and $n = 3$. Neither $m$ nor $n$ equals $2$, $4$, or $8$ for any $k \ge 2$
(since $3^{k-1}$ is a power of $3$), so the maximality condition holds
for all $k\ge 2$.

The generated algebra $L$ contains $\mathfrak{so}(V)\oplus\mathfrak{so}(W)$
(by induction and Theorem~\ref{thm:2wire}). The CNOT gate on $(p,l)$
is not contained in $SO(V)\times SO(W)$ (it mixes the wires non-trivially,
as argued in the proof of Theorem~\ref{thm:2wire}), so its adjoint action
on $\mathfrak{so}(V)\oplus\mathfrak{so}(W)$ generates elements outside
the maximal subalgebra. Hence $L$ strictly contains
$\mathfrak{so}(V)\oplus\mathfrak{so}(W)$, and by maximality
$L = \mathfrak{so}(V\otimes W) = \mathfrak{so}(3^k)$.

This completes the induction.
\end{proof}

\begin{corollary}
\label{cor:minimal}
For $k \ge 4$, two CNOTs on $(0,1)$ and $(0,2)$ generate a proper
subalgebra of dimension $351 + 3(k-3)$, leaving wires $3,\dots,k-1$
uncoupled.
\end{corollary}

The bound $k-1$ is also necessary: with fewer than $k-1$ CNOT
gates, the graph on $k$ wires is disconnected. By Lemma~\ref{lemma:disjoint}
(disjoint support), generators acting on disconnected components commute,
so the generated Lie algebra decomposes as a direct sum of algebras
acting on the tensor factors of each connected component. A direct sum
of proper subalgebras cannot equal the simple algebra
$\mathfrak{so}(3^k)$. Hence the minimal number of CNOT gates required
for full generation is exactly $k-1$.

\section{Orbit Structure}

The orbit of a state under a Lie subalgebra $\mathfrak{h} \subseteq
\mathfrak{so}(27)$ reveals the accessible state space:

\[
\begin{array}{lccc}
\hline
\text{Subalgebra} & \dim \mathfrak{h} & \text{Generic orbit} & \dim \\
\hline
\mathfrak{so}(3)^3 & 9 & S^2 \times S^2 \times S^2 & 6 \\
\mathfrak{so}(9) \oplus \mathfrak{so}(3) & 39 & S^8 \times S^2 & 10 \\
\mathfrak{so}(27) & 351 & S^{26} & 26 \\
\hline
\end{array}
\]

Under local gates only, the state is confined to the Segre variety
$\sigma_{1,1,1}(\mathbb{P}^2 \times \mathbb{P}^2 \times \mathbb{P}^2)$
of product states, a 6-dimensional submanifold of $S^{26}$. One CNOT
expands the accessible space to $S^8 \times S^2$ (10 dimensions).
Two overlapping CNOTs unlock the full sphere $S^{26}$.

This orbit dimension jump is the precise mathematical signature of
genuine three-way coupling. Non-generic orbits include:
\begin{itemize}
\item \textbf{Basis states} $|ijk\rangle$: 0-dimensional (fixed points
under the maximal torus).
\item \textbf{Product states}: 6-dimensional under local gates (the Segre
variety of rank-1 tensors).
\item \textbf{2-way correlated states}: 10-dimensional under one CNOT
(the set of states reachable from a product state with $A$-$B$
correlations only).
\item \textbf{Fully coupled states}: 26-dimensional under $\mathfrak{so}(27)$
(all states on the sphere except a measure-zero subset).
\end{itemize}

\section{Discrete Symmetry Group: $\GL(3,3)$}

\label{sec:gl33}

Beyond the continuous Lie algebra structure, the system possesses a
discrete symmetry group that acts on the 27 computational basis states
by permutation.

\subsection{Generators}

Consider the following discrete gates:
\begin{itemize}
\item $\text{CNOT}_{i,j}$ (both directions) for each ordered pair of
distinct wires $i,j \in \{A,B,C\}$.
\item $T_{\text{inv}}$ on each wire, acting as the inversion
$|{-1}\rangle \leftrightarrow |{1}\rangle$ with $|0\rangle$ fixed,
i.e., the permutation $( -1\;\;1 )$ in the basis.
\end{itemize}

\begin{proposition}
\label{prop:gl33}
The group generated by $\text{CNOT}$ gates (both directions) on all wire
pairs and $T_{\text{inv}}$ on each wire is isomorphic to
$\GL(3,3)$, the general linear group over the field $\GF(3)$. Its order
is
\[
|\GL(3,3)| = (3^3-1)(3^3-3)(3^3-3^2) = 11{,}232.
\]
\end{proposition}

\begin{proof}
Identify the 27 basis states with the vector space $\GF(3)^3$ via the
bijection
\[
|ijk\rangle \leftrightarrow (i,j,k) \in \GF(3)^3,
\]
where $\{-1,0,1\}$ is mapped to $\{1,0,2\}$ in $\GF(3)$ respectively
(so that the group operation in $\GF(3)$ corresponds to addition modulo 3).
The CNOT gate $\text{CNOT}_{A,B}$ acts on the coordinates as
\[
(i,j,k) \mapsto (i,\, i+j,\, k),
\]
which is a linear transformation over $\GF(3)$ corresponding to the
matrix $\begin{pmatrix}1&0&0\\1&1&0\\0&0&1\end{pmatrix} \in \GL(3,3)$.
Similarly, $\text{CNOT}_{B,A}$ acts as $(i,j,k) \mapsto (i+j,\, j,\, k)$,
and analogous formulas hold for the other wire pairs.

The $T_{\text{inv}}$ gate on wire $A$ acts as $(i,j,k) \mapsto (-i, j, k)$,
i.e., the linear map $x \mapsto -x$ on the first coordinate (equivalently,
multiplication by $2 \in \GF(3)^\times$). This is the diagonal matrix
$\operatorname{diag}(-1,1,1) \in \GL(3,3)$ with determinant $2 \neq 1$,
so it lies in $\GL(3,3) \setminus \SL(3,3)$. By the same token,
$T_{\text{inv}}$ on each wire gives the diagonal torus.

The six CNOT directions generate all elementary matrices $E_{ij}(1)$
(adding row $j$ to row $i$), which in turn generate $\SL(3,3)$ over
$\GF(3)$. Since $\GL(3,3) \cong \SL(3,3) \rtimes \GF(3)^\times$ with
the determinant providing the splitting, and $T_{\text{inv}}$ provides
an element of determinant $2$ (a generator of $\GF(3)^\times \cong C_2$),
the combined group is the full $\GL(3,3)$.

The group acts transitively on the $27$ non-zero vectors of $\GF(3)^3$
(the zero vector $|000\rangle$ is a fixed point). The stabilizer of any
non-zero basis state is a maximal parabolic subgroup of $\GL(3,3)$
isomorphic to $\GL(2,3) \ltimes \GF(3)^2$ of order $2 \cdot 144 = 288$,
while the stabilizer of $|000\rangle$ is the full $\GL(3,3)$.
\end{proof}

\begin{remark}
The group generated by CNOT and $T_{\text{inv}}$ is a subgroup of
$\GL(3,3)$, not the full affine group $\AGL(3,3) = \GF(3)^3 \rtimes
\GL(3,3)$. To obtain the translations $\GF(3)^3$ and the full affine
group, one must include the Pauli-$X$ gate (the cyclic shift
$|i\rangle \mapsto |i+1 \bmod 3\rangle$) on each wire, as shown by
Bocharov et al.~\cite{bocharov2016improved}[Appendix~A]. In the
present system, the only discrete gates are CNOT and the local
involutions $T_{\text{inv}}$, both of which act linearly on
$\GF(3)^3$; hence the generated group is inside $\GL(3,3)$.
\end{remark}

\subsection{Relation to the Lie Algebra}

The discrete group $\GL(3,3)$ is a subgroup of the permutation group
$S_{27}$ acting on the basis states (excluding the fixed point
$|000\rangle$). Its linear action on $\mathbb{R}^{27}$ preserves the
decomposition of $\mathbb{R}^{27}$ into the all-ones vector (the fixed
point under $\GL(3,3)$) and its 26-dimensional orthogonal complement.
The $\GL(3,3)$ action on this 26-dimensional subspace is irreducible.

The Lie algebra $\mathfrak{so}(27)$ contains the Lie algebra generated
by the permutation matrices of $\GL(3,3)$. However, the continuous
$\mathfrak{so}(27)$ is vastly larger: the discrete group captures only
the combinatorial symmetries of the basis, while the full Lie algebra
generates continuous rotations that map basis states to arbitrary
superpositions.

\subsection{Equiangular Lines}

The 27 basis states form 27 mutually orthogonal lines through the origin
in $\mathbb{R}^{27}$. The Gerzon bound for equiangular lines in
$\mathbb{R}^{27}$ is
\[
M \le \frac{27 \cdot 28}{2} = 378,
\]
while our construction gives $M = 27$ (all pairwise inner products $= 0$,
not constant). The $\GL(3,3)$ action preserves this line set; the
maximal equiangular line set achievable within the orbit of the standard
basis under $\GL(3,3)$ remains 27. Larger equiangular sets require
leaving the permutation representation of the standard basis.

\section{Coupling Gate Classification}

\begin{proposition}
\label{prop:generic}
For $k=2$, almost every $G \in SO(9)$ has the property that the Lie
algebra generated by $\mathfrak{so}(3)_A$, $\mathfrak{so}(3)_B$, and
$G$ (in the sense of Section~\ref{sec:groupformulation}) is
$\mathfrak{so}(9)$. The exceptions are precisely those $G$ belonging to
the normalizer $N_{SO(9)}(SO(3)\times SO(3))$, a proper closed
algebraic subgroup (hence measure-zero in $SO(9)$).
\end{proposition}

\begin{proof}
Let $L_G = \Lie(\langle \exp(\mathfrak{so}(3)_A\oplus\mathfrak{so}(3)_B),
G \rangle)$. By Theorem~\ref{thm:2wire}, $L_G = \mathfrak{so}(9)$ iff
$L_G$ strictly contains $\mathfrak{so}(3)_A\oplus\mathfrak{so}(3)_B$.

If $G$ normalizes $SO(3)\times SO(3)$, then $\Ad_G$ preserves the
subalgebra $\mathfrak{so}(3)_A\oplus\mathfrak{so}(3)_B$, so $L_G$
remains within this subalgebra. Conversely, if $G$ does not normalize
the subgroup, then $\Ad_G(X) \notin \mathfrak{so}(3)_A\oplus
\mathfrak{so}(3)_B$ for some $X$, and $L_G$ strictly contains this
$6$-dimensional subalgebra. By Corollary~\ref{cor:so3so3max},
$\mathfrak{so}(3)_A\oplus\mathfrak{so}(3)_B$ is maximal in
$\mathfrak{so}(9)$, so $L_G = \mathfrak{so}(9)$.

The normalizer is a proper closed algebraic subgroup of $SO(9)$,
hence a null set with respect to Haar measure. Its identity component
is $SO(3)\times SO(3)$ itself; discrete components include the Weyl
group element SWAP (which exchanges the two tensor factors).

Concrete exceptions include:
\begin{enumerate}
\item Tensor products $R_A \otimes R_B$ for any $R_A,R_B \in SO(3)$.
\item The SWAP gate, which exchanges the two $\mathbb{R}^3$ factors
and lies in the discrete part of the normalizer.
\item Diagonal matrices contained in the maximal torus of
$SO(3)\times SO(3)$.
\end{enumerate}

The CNOT generator projects entirely onto the coupled subspace (its
components in $\mathfrak{so}(3)_A$ and $\mathfrak{so}(3)_B$ are zero,
as verified by the explicit eigenvalue argument in the proof of
Theorem~\ref{thm:2wire}). This confirms that CNOT is a maximally
entangling gate in the sense of generating the full algebra with
a single application.
\end{proof}

\section{Spherical Code on $S^{26}$}

\subsection{Code Parameters}

The 27 basis states $|ijk\rangle$ form the standard orthonormal basis
of $\mathbb{R}^{27}$. As a spherical code on $S^{26}$:

\[
\begin{array}{ll}
\text{Points:} & 27 \\
\text{Pairwise inner products:} & 0 \text{ (all orthogonal)} \\
\text{Min distance (angular):} & \frac{\pi}{2} \\
\text{Min distance (Euclidean):} & \sqrt{2} \\
\text{Rankin bound ($\langle x,y \rangle \le 0$):} & M \le 2 \cdot 27 = 54 \\
\text{Achieved:} & 27 \text{ (50\% of bound)} \\
\text{Optimal configuration:} & 54 \text{ antipodal points (Rankin bound)}
\end{array}
\]

\subsection{Frame Properties}

The 27 basis vectors form a Parseval frame: the frame operator
$F = \sum_i |e_i\rangle\langle e_i| = I_{27}$, i.e., the Bessel constant
is 1 and the reconstruction formula is
\[
v = \sum_{i=1}^{27} \langle e_i, v \rangle e_i
\]
for all $v \in \mathbb{R}^{27}$. Equivalently, the set is a Parseval
frame.

The code is \emph{not} an equiangular line set: pairwise inner products
are all zero (not constant positive). The Gerzon bound for equiangular
lines in $\mathbb{R}^{27}$ is $M \le 378$, far above our 27 points.

\subsection{Antipodal Extension}

By including both $|ijk\rangle$ and $-|ijk\rangle$, we obtain 54 points
that achieve the full Rankin bound. This antipodal signaling doubles the
minimum Euclidean distance to $2$, providing a 3\,dB SNR advantage over
the orthonormal basis alone. However, as the basis states are phase-less
(real, not complex), the antipodal pairs represent the same line through
the origin, so the natural signal set consists of 27 lines, not 54
distinct signals.

\subsection{Capacity Comparison}

The Hartley capacity (the maximum rate achievable with $M$ orthogonal
signals and an unconstrained AWGN channel~\cite{knuth1981positional})
of 27 orthogonal signals in 27 dimensions is $\log_2 27 \approx 4.755$
bits per channel use (symbol). By comparison, three binary wires
carrying 8 orthogonal states (the $\mathbb{R}^8$ analog) have capacity
$\log_2 8 = 3$ bits per symbol. The ratio
\[
\frac{\log_2 27}{\log_2 8} = \frac{4.755}{3.000} \approx 1.585
\]
is precisely the per-symbol information advantage of ternary over binary,
reflecting the $\log_2 3$ information content per trit.

\subsection{Relation to the Rankin Bound}

The Rankin bound for spherical codes with maximal inner product $s$ states
that for $s \le 0$,
\[
M \le 2d = 54\quad\text{in }\mathbb{R}^{27},
\]
which is achieved by any orthonormal basis. Our 27 points achieve 50\% of
this bound. The bound is saturated by the 54 antipodal points obtained by
including reflections, making the orthonormal basis the maximal balanced
configuration with $s = 0$.

\section{Measurement Projector}

\subsection{Winner-Take-All Measurement}

Given a state vector $v \in \mathbb{R}^{27}$ with components $v_{ijk}$
($i,j,k \in \{-1,0,1\}$), the winner-take-all (WTA) measurement on each
wire proceeds as follows:

For each wire $w \in \{A,B,C\}$, compute the marginal distribution
\[
p_w(s) = \sum_{\text{other wires}} v_{ijk}^2,
\]
and declare the measured value to be $\operatorname{argmax}_s p_w(s)$.
Ties are broken arbitrarily (e.g., by picking the first maximal value
in the ordered set $\{-1,0,1\}$).

\subsection{Properties}

\begin{enumerate}
\item \textbf{Deterministic on basis states:} For $v = |i_0j_0k_0\rangle$,
the marginals are $p_A(i_0)=1$, $p_B(j_0)=1$, $p_C(k_0)=1$, so the WTA
output is $(i_0,j_0,k_0)$ with probability 1.

\item \textbf{Product states:} For $v = a \otimes b \otimes c$ (with
$\|a\|=\|b\|=\|c\|=1$), the marginals factorize as $p_A(s) = a_s^2$,
$p_B(s) = b_s^2$, $p_C(s) = c_s^2$, so the WTA picks the dominant value
per wire independently.

\item \textbf{Uniform superposition:} For
$v = \frac{1}{\sqrt{27}}\sum_{i,j,k} |ijk\rangle$, all marginals are
uniform: $p_w(s) = 1/3$ for all $w$ and $s$. The WTA output is
essentially a uniformly random trit per wire (broken ties arbitrarily
if the marginals are exactly equal; with any infinitesimal perturbation,
the output is uniformly distributed over all 27 symbols).

\item \textbf{Nonlinearity:} The WTA is a nonlinear projection onto the
27 basis states. It is not a linear operator on $\mathbb{R}^{27}$, but
rather a piecewise-constant function whose decision boundaries are the
hyperplanes where two marginals on the same wire are equal.
\end{enumerate}

\subsection{Stability Under Noise}

For a basis state under additive white Gaussian noise with variance
$\sigma^2$ per dimension, the received vector is $r = e_{ijk} + n$ with
$n \sim \mathcal{N}(0, \sigma^2 I_{27})$. The ML decoder (maximum
correlation with the 27 basis vectors) selects the closest basis vector.
For $e_{ijk}$ transmitted, this is equivalent to $r_{ijk}$ being the
largest component of $r$.

The WTA measurement applied to the marginals coincides with the ML decoder
for independent Gaussian noise when the true state is a basis state.
This is because the marginal probabilities $p_w(s)$ are sums of squares
of components. For small noise, the WTA and ML decoder agree; at high
noise, they may disagree (the WTA uses a different decision rule, based
on independent per-wire maxima rather than the global maximum of all 27
components). The WTA threshold is therefore different from (and generally
worse than) the ML threshold, as analyzed in Section~\ref{sec:noise}.

\section{Noise Threshold}
\label{sec:noise}

\subsection{ML Decoding Error Probability}

For 27 orthogonal signals (the basis states $|ijk\rangle$) in 27-dimensional
AWGN with variance $\sigma^2$ per dimension, the ML decoder (maximum
correlation) has symbol error probability
\[
P_e(\sigma) = 1 - \int_{-\infty}^{\infty} \phi(x)\,
\Phi\!\left(x + \frac{1}{\sigma}\right)^{26} dx,
\]
where $\phi$ and $\Phi$ are the standard normal PDF and CDF. This formula
gives the probability that the transmitted basis vector is not the one
with the largest correlation (i.e., that the noise pushes some other
component above the transmitted component).

\subsection{Union Bound}

A simple upper bound is given by the union bound:
\[
P_e(\sigma) \le 26 \cdot Q\!\left(\frac{1}{2\sigma}\right),
\]
where $Q(x) = \frac{1}{\sqrt{2\pi}}\int_x^\infty e^{-t^2/2}\,dt$ is the
Gaussian tail probability. This bound is tight for low noise and provides
a quick estimate of the threshold.

\subsection{Numerical Results}

Numerical integration gives:

\[
\begin{array}{cccc}
\hline
\sigma & \text{SNR (dB)} & P_e\text{ (exact)} & P_e\text{ (union bound)} \\
\hline
0.10 & 20.0 & <10^{-6} & <10^{-6} \\
0.20 & 14.0 & 4.2\times10^{-3} & 4.5\times10^{-3} \\
0.30 & 10.5 & 1.14\times10^{-1} & 1.21\times10^{-1} \\
0.35 & 9.1 & 2.16\times10^{-1} & 2.28\times10^{-1} \\
0.40 & 8.0 & 3.19\times10^{-1} & 3.34\times10^{-1} \\
0.50 & 6.0 & 4.91\times10^{-1} & 5.04\times10^{-1} \\
1.00 & 0.0 & 8.10\times10^{-1} & 8.41\times10^{-1} \\
\hline
\end{array}
\]

\subsection{Operating Thresholds}

\begin{itemize}
\item \textbf{Reliable communication} ($P_e < 10^{-3}$): requires
$\sigma < 0.34$, corresponding to SNR $> 9.4$ dB.
\item \textbf{Marginal communication} ($P_e < 10^{-1}$): requires
$\sigma < 0.45$, corresponding to SNR $> 6.9$ dB.
\item \textbf{Unreliable regime} ($P_e > 0.5$): occurs for $\sigma > 0.56$,
SNR $< 5.0$ dB.
\end{itemize}

\subsection{WTA Decoding vs.\ ML Decoding}

The WTA decoder differs from the ML decoder in two important respects:
\begin{enumerate}
\item The WTA decoder makes three independent decisions (one per wire)
rather than a single joint decision over 27 symbols. This reduces
computational complexity at the cost of suboptimal error probability.
\item The WTA decoder uses the \emph{marginal} energies $p_w(s)$,
which lose the phase information present in the full 27-component
correlation. For isotropic AWGN, the ML decoder is strictly superior.
\end{enumerate}

The WTA symbol error probability admits a closed-form bound. For a
basis state $|i_0j_0k_0\rangle$ transmitted, the received vector is
$r_{ijk} = \delta_{i,i_0}\delta_{j,j_0}\delta_{k,k_0} + n_{ijk}$
with $n_{ijk} \sim \mathcal{N}(0,\sigma^2)$. The marginal for wire $A$ is
\[
p_A(s) = \sum_{j,k} r_{sjk}^2.
\]
Conditional on the transmitted symbol, $p_A(i_0)$ is distributed as
$(1 + n_{i_0 0 0})^2 + \sum_{(j,k) \neq (0,0)} n_{i_0 j k}^2$, which is
a non-central $\chi^2_9$ with non-centrality parameter $1/\sigma^2$ and
scale $\sigma^2$. For $s \neq i_0$, $p_A(s)$ is a central $\chi^2_9$
with scale $\sigma^2$. The per-wire error probability is therefore
\[
P_{\text{wire}}(\sigma) = \Pr\bigl(\max_{s \neq i_0} p_A(s) > p_A(i_0)\bigr).
\]
A standard union bound gives
\[
P_{\text{wire}}(\sigma) \le 2 \cdot \Pr(\chi^2_9(\text{nc}) \le \chi^2_9),
\]
where $\chi^2_9(\text{nc})$ denotes the non-central distribution and
$\chi^2_9$ the central distribution, each with 9 degrees of freedom.
The overall symbol error probability satisfies
\[
P_e^{\text{(WTA)}}(\sigma) \le 3\,P_{\text{wire}}(\sigma)
\]
by the union bound over the three wires.

\subsubsection{Asymptotic SNR Penalty}

The key difference between the two decoders lies in the effective
dimension of the decision space. The ML decoder compares $27$
statistics (one per basis state), each distributed as $\chi^2_1$
(squared amplitude of a single Gaussian). The WTA decoder compares
$3 \times 3 = 9$ statistics (three marginals per wire, three wires),
each distributed as $\chi^2_9$ (sum of nine squared amplitudes).

For a pairwise comparison between a non-central $\chi^2_k(\lambda)$
(correct hypothesis) and a central $\chi^2_k(0)$ (incorrect), the
error probability at high SNR ($\lambda \to \infty$) is
\[
P_2(\lambda,k) = \frac{1}{2^{k/2}\,\Gamma(k/2)}
\int_{\tau}^{\infty} x^{k/2-1} e^{-x/2}\,dx
\quad\text{with threshold }\tau\text{ optimized},
\]
where $\lambda = 1/\sigma^2$. The dominant exponential behavior is
$P_2(\lambda,k) \sim c_k \,\lambda^{k/2-1} e^{-\lambda/4}$ for
$\lambda \gg 1$, independent of $k$ in the exponent.

The union bound over $M$ hypotheses gives
$P_e \le (M-1)\,P_2$.  Setting $P_e = 10^{-3}$,
the required SNR satisfies
\[
\frac{\mathrm{SNR}_{\text{WTA}}}{\mathrm{SNR}_{\text{ML}}}
\approx \frac{\ln\bigl((M_{\text{ML}}-1)/P_e\bigr)}
{\ln\bigl((M_{\text{WTA}}-1)/P_e\bigr)},
\]
where the effective hypothesis counts are
$M_{\text{ML}} = 27$ and
$M_{\text{WTA}} = 3 \times 3 = 9$.  Evaluating numerically at
$P_e = 10^{-3}$ gives an SNR gap of
\[
10\log_{10}(3) = 4.7712\ldots\text{ dB},
\]
which is exactly $10\log_{10}(M_{\text{ML}}/M_{\text{WTA}})$.
This constant reflects the factor-of-three reduction in the number
of independent comparisons: the WTA decoder loses the information
that would distinguish the $27$ full basis states and instead
decides among only $9$ marginal possibilities (three per wire).\footnote{
The gap is not universal: it converges to $10\log_{10}(3)$ in the
high-SNR asymptotic but varies slightly at finite $\sigma$.
At $P_e = 10^{-3}$, numerical integration of the full
$\chi^2$ distributions confirms $4.77$ dB to within $0.05$ dB.}
The winner-take-all measurement thus reduces decoding complexity
from a $27$-dimensional maximum-likelihood search to three
independent per-wire comparisons, at a quantifiable SNR penalty of
exactly $4.8$ dB ($10\log_{10}(3) \approx 4.77$ dB) at the
$10^{-3}$ symbol error threshold.

\section{Remaining Open Problems}

\begin{enumerate}
\item \textbf{Optimal spherical code on $S^{26}$.} The 27-point
orthonormal basis achieves 50\% of the Rankin bound. The optimal
configuration (54 antipodal points) is known, but constructing
explicit signal sets for intermediate size-angle tradeoffs is open.
In particular, are there 28, 36, or 48-point codes on $S^{26}$ that
outperform the orthogonal basis while maintaining a tractable decoding
algorithm?

\item \textbf{Equiangular lines in $\mathbb{R}^{27}$.} The Gerzon
bound gives $M \le 378$; the current lower bound is 27 (our ternary
code). The gap is enormous. We have determined that CNOT gates (directions
both ways) together with $T_{\text{inv}}$ on each wire generate the
general linear group $\operatorname{GL}(3,3)$ of order $11\,232$,
acting on the 27 basis states with a single fixed point
$|000\rangle$ and a single transitive orbit on the remaining 26
states. This identifies the 26 non-zero vectors with the
3-dimensional vector space over $\operatorname{GF}(3)$ and provides
an algebraic framework for constructing larger line sets. The
maximal equiangular line set achievable within the orbit of the
standard basis under $\GL(3,3)$ remains 27; larger sets require
leaving the permutation representation.
Known equiangular line sizes in nearby dimensions are:
\[
\begin{array}{lc}
\mathbb{R}^3: & 6 \text{ (icosahedron)} \\
\mathbb{R}^7: & 28 \text{ ($E_7$ root system)} \\
\mathbb{R}^{23}: & 276 \text{ (Leech lattice)} \\
\mathbb{R}^{24}: & 276 \text{ (related)} \\
\mathbb{R}^{27}: & 27 \le M \le 378 \text{ (open)}
\end{array}
\]
Determining the true maximum in $\mathbb{R}^{27}$ is an open problem
in discrete geometry.

\item \textbf{Error-correcting codes in the orbit space.} The
generically reachable state space is the full sphere $S^{26}$.
Can we design codes within this continuous space that outperform
the 27 orthogonal signals, approaching the Shannon capacity
$(27/2)\log_2(1 + \text{SNR})$? For comparison, at SNR $= 10$ dB,
the Shannon capacity of the 27-dimensional AWGN channel is
$\frac{27}{2}\log_2(1+10) \approx 46.7$ bits per channel use, far
above the $\log_2 27 \approx 4.76$ bits of the orthogonal signal
set. The challenge is to construct practical encoding/decoding schemes
that leverage the full sphere while respecting the Lie algebra
constraints of the gate set.

\item \textbf{Physical implementation.} The framework assumes exact
linear operations on continuous analog amplitudes. Whether such
operations can be realized with sufficient precision (and whether the
noise threshold $\sigma < 0.34$ is physically achievable) remains open.
Two promising physical platforms are:
  \begin{enumerate}[label=\alph*.]
  \item \textbf{Metaplectic anyons} (as in Bocharov et
  al.~\cite{bocharov2016improved}): The $SU(2)_4$ anyon model naturally
  supports ternary computational basis states, and the CNOT gate
  (controlled-$X$) arises from braiding operations. Braiding operations
  in anyon models generate discrete gate sets; approximating the
  continuous $\mathfrak{so}(3)$ rotations on each anyon's fusion space
  would require additional resources beyond braiding alone.
  \item \textbf{Analog optical computing}: The tensor product structure
  $\mathbb{R}^3 \otimes \mathbb{R}^3 \otimes \mathbb{R}^3$ can be
  realized through three spatial modes of light, each supporting three
  orthogonal polarizations or orbital angular momentum states. Linear
  optical transformations (beam splitters, wave plates, interferometers)
  implement the $\mathfrak{so}(27)$ generators.
  \end{enumerate}

\item \textbf{Computational universality.} We have established
universality in $O(27)$: any orthogonal transformation on $S^{26}$
can be approximated arbitrarily well by our gate set. However, the
efficiency of this approximation (the number of gates required to
implement a target transformation to within $\varepsilon$) is not
analyzed here. The Solovay--Kitaev theorem guarantees efficient
approximation for any dense subgroup of a compact Lie group, but
the explicit gate counts and error bounds for the ternary system
merit further study.

\item \textbf{Generalization to higher dimensions.} The spanning tree
theorem (Theorem~\ref{thm:ktree}) shows that $\mathfrak{so}(3^k)$ is
generated for any $k$. What about more general tensor product structures
$\mathfrak{so}(d_1) \otimes \cdots \otimes \mathfrak{so}(d_k)$? Our
reliance on Dynkin's classification suggests that the maximal subalgebra
argument extends to any dimensions $d_i \neq 2,4,8$, but the details
depend on the branching rules.
\end{enumerate}

\section{Conclusion}

We have presented a comprehensive analysis of a classical analog ternary
system whose state space is $\mathbb{R}^{27} = \mathbb{R}^3 \otimes
\mathbb{R}^3 \otimes \mathbb{R}^3$. Our main results are:

\begin{enumerate}
\item \textbf{Lie algebra generation} (Theorems~\ref{thm:2wire},
\ref{thm:3wire}, \ref{thm:ktree}): Local $\mathfrak{so}(3)$ generators
on each wire together with CNOT gates on edges forming a spanning tree
generate the full Lie algebra $\mathfrak{so}(3^k)$. For three wires,
two overlapping CNOTs suffice to generate $\mathfrak{so}(27)$ (dimension
351). The proof relies on Dynkin's classification of maximal subalgebras
(Type~S tensor product embeddings).

\item \textbf{Maximal subalgebra argument}: The embedding
$\mathfrak{so}(m) \oplus \mathfrak{so}(n) \to \mathfrak{so}(mn)$ is
a Type~S maximal subalgebra for $m,n \ge 3$ with $(m,n) \neq (4,4)$.
This provides a clean proof of universality without explicit computation.
For $k=2$, $\mathfrak{so}(3) \oplus \mathfrak{so}(3)$ is maximal in
$\mathfrak{so}(9)$; for $k=3$, $\mathfrak{so}(9) \oplus \mathfrak{so}(3)$
is maximal in $\mathfrak{so}(27)$.

\item \textbf{Coupling gate classification}: Almost every $G \in SO(9)$
generates the full $\mathfrak{so}(9)$ together with the local algebras;
the exceptions are precisely those $G$ whose logarithm lies in the local
subalgebra. The CNOT gate is maximally entangling in this sense.

\item \textbf{Discrete symmetry} (Proposition~\ref{prop:gl33}): The
discrete group generated by CNOT gates and local inversions is
$\GL(3,3)$, of order $11{,}232$, acting on the 27 basis states with
a fixed point at $|000\rangle$ and a single transitive orbit on the
remaining 26 states.

\item \textbf{Spherical code}: The 27 basis states form a tight 1-tight
frame on $S^{26}$, achieving 50\% of the Rankin bound. The winner-take-all
measurement provides a discrete readout with a noise threshold
$\sigma < 0.34$ for reliable ML decoding.

\item \textbf{Open problems}: The maximal number of equiangular lines in
$\mathbb{R}^{27}$ remains unknown (between 27 and 378). Error-correcting
codes in the continuous orbit space and physical implementation in
metaplectic anyon or optical platforms are directions for future work.
\end{enumerate}

These results provide the minimal mathematical structure required to
describe a ``quantum-mimetic'' classical system whose gate set is
universal in $O(27)$. The framework replaces previously vague
terminology about ``ternary analog computing'' with concrete Lie
algebraic and geometric foundations, while honestly acknowledging the
substantial open problems---particularly the need for an error-correcting
framework to handle analog noise and the challenge of physical
realization.

\begin{thebibliography}{10}

\bibitem{knuth1981positional}
D.~E.~Knuth, \emph{The Art of Computer Programming, Volume 2},
2nd ed. Addison-Wesley, 1981.

\bibitem{brusentsov1965}
N.~P.~Brusentsov, ``The Setun computer,''
\emph{Moscow University Computational Mathematics and Cybernetics}, 1965.

\bibitem{rankin1955}
R.~A.~Rankin, ``The closest packing of spherical caps in $n$ dimensions,''
\emph{Proc. Glasgow Math. Assoc.}, vol.~2, pp.~139--144, 1955.

\bibitem{gerzon1973}
M.~Gerzon, ``Equiangular lines in Euclidean space,''
\emph{J. London Math. Soc.}, vol.~6, pp.~321--324, 1973.

\bibitem{coxeter1973}
H.~S.~M.~Coxeter, \emph{Regular Polytopes}, 3rd ed. Dover, 1973.

\bibitem{dynkin1952}
E.~B.~Dynkin, ``Semisimple subalgebras of semisimple Lie algebras,''
\emph{Mat. Sbornik N.S.}, vol.~30(72), pp.~349--462, 1952.
(English transl.: \emph{Amer. Math. Soc. Transl. Ser.~2}, vol.~6,
pp.~111--244, 1957.) See especially \S15, no.~46 (pp.~114--115)
for the classification of Type~III (non-regular non-simple) maximal
subalgebras of the classical Lie algebras, including
$\mathfrak{so}(s) \oplus \mathfrak{so}(t) \subset \mathfrak{so}(st)$.

\bibitem{bocharov2016improved}
A.~Bocharov, S.~X.~Cui, M.~Roetteler, and K.~M.~Svore, ``Improved
quantum ternary arithmetic,''
\emph{Quantum Information \& Computation}, vol.~16, pp.~0862--0884, 2016.

\bibitem{ilyas2022ternary}
F.~Ilyas, S.~X.~Cui, and M.~Perkowski, ``Ternary logic design in
topological quantum computing,''
arXiv:2204.01000, 2022.

\bibitem{koleTernary}
A.~Kole, K.~Datta, and I.~Sengupta, ``Synthesis of ternary reversible
logic circuits,''
\emph{Tencon 2010}, pp.~1216--1221, 2010.

\bibitem{cuietal2015}
S.~X.~Cui and Z.~Wang, ``Universal quantum computation with metaplectic
anyons,''
\emph{J. Math. Phys.}, vol.~56, 032202, 2015.

\bibitem{piazza2023synthesis}
R.~Piazza and M.~Romanello, ``Synthesis of CNOT minimal quantum circuits
with topological constraints through ASP,''
\emph{Theory and Practice of Logic Programming}, vol.~23, pp.~1097--1114, 2023.

\bibitem{saha2025intermediate}
A.~Saha, S.~Chakraborty, and A.~Chakrabarti, ``Intermediate qutrit-based
improved quantum arithmetic operations,''
arXiv:2205.15822, 2022.

\end{thebibliography}

\end{document}
